\scanpage{65}
\mspage{II. 13}
\uline{Fin de la démonstration de} 1.9.

\uline{Bonne réduction} \deleted{\unsure{aux pts de car. $\ell$}} (alias cas \unsure{nilpotent|unipotent}).

% lines 3-4
a) L'action est semi-simple et abélienne.

On utilise les faits suivants :

% lines 5-7
\uline{Lemme} Si $A$ est un schéma abélien sur un corps fini $k$, l'action du frobénius \inserted{géométrique} $f$ sur $T_\ell(A(\overline k))$ est semi-simple.

% lines 8-15
En effet, \deleted{\unsure{or}} \deleted{\unsure{c'est un}} \inserted{cette action est \deleted{précisément} \inserted{inverse de} celle du frobénius géométrique} \uline{endomorphisme} de $A_{\unsure{\overline k|k|\xi}}$. \deleted{Donc} Or cet $f$ appartient au \uline{centre} de l'anneau $\operatorname{End}(A)\otimes_{\mathbb{Z}}\mathbb{Q}$ des endomorphismes, \inserted{tensorisé par $\mathbb{Q}$,} qui, par Weil [2] est une algèbre \inserted{finie} semi-simple sur \deleted{\unsure{A}} $\mathbb{Q}$, donc son centre $Z$ est une \deleted{\unsure{extension}} algèbre finie \deleted{\unsure{séparable}} \inserted{étale} de $\mathbb{Q}$. \inserted{Donc $Z_\ell=Z\otimes_{\mathbb{Q}}\mathbb{Q}_\ell$ \deleted{\unsure{opère}} est une algèbre finie étale sur $\mathbb{Q}_\ell$;} \deleted{Donc} \deleted{$Z$ opère} \unsure{alors} il est clair $Z$ opère de façon semi-simple sur $T_\ell(A)\otimes_{\mathbb{Z}_\ell}\mathbb{Q}_\ell$,

% lines 16-24
\deleted{\uline{Lemme}} \inserted{\uline{Corollaire} 1} Soit donc $S$ un schéma \inserted{loc} de t.f. sur $\operatorname{Spec}\mathbb{Z}$, $\unsure{\xi|\zeta}$ un pt géom. de $S$, $A$ un schéma abélien sur $S$, \inserted{$\ell$ un nb premier distinct des car. résiduelles de $S$,} d'où une représentation de $\pi_1(S,\xi)$ sur \inserted{$E=$} $T_\ell(A(\xi))\otimes_{\mathbb{Z}_\ell}\mathbb{Q}_\ell$.

Pour tout point fermé $s\in S$, les éléments de frobénius $f_s\in\pi_1(S,\xi)$ [déterminé à \deleted{\unsure{auto}} conj. près comme l'image du gén. de frob. de $\pi_1(s)\deleted{\unsure{\xi}}$ par $\pi_1(s)\mathrel{\text{---}}\pi_1(S)$] opère \deleted{sur} de façon semi-simple.

% lines 25-27
\uline{Corollaire} 2 Si la représentation de $\pi_1(S,\xi)$ sur $E$ est \deleted{\unsure{si}} ess\textsuperscript{t} nilpotente, elle est \deleted{abélienne et} semi-simple et ess. abélienne.
% archive sheet number: 64
